% ======================================================================
% Main-text version of the two-solution theory. Full proofs of
% Lemmas ~\ref{lem:clones}, \ref{lem:decode}, \ref{lem:ema},
% Theorem~\ref{thm:lambdastar}, Proposition~\ref{prop:kn}, and
% Corollary~\ref{cor:bipartite}, together with the measurement
% protocol, census, and audits, are in Appendix~\ref{app:tworoutes}.
% ======================================================================
\section{The Two-Solution Structure of Min-sum on SCFGs}\label{sec:tworoutes}

We now explain the behavior documented in Section~\ref{Sec:MS-split} in
three conditional steps: (i) a split function-node whose incoming
difference is in a regime transmits a \emph{decision} --- one row of its
cost table --- rather than a magnitude, and when this holds for every
message involved (\emph{full commitment}), the selected assignments evolve
by a simple synchronous rule; (ii) \emph{any} assignment sequence
following that rule is eventually periodic with period one or two, and the
limiting pair is a local minimum of a pair cost we call the alternation
cost; (iii) damping destroys a fixed committed period-2 pattern exactly
above an explicit threshold, and the threshold approaches $1$ as the
degree grows, so no single damping value suffices universally
(Section~\ref{sec:whydamping}).  Throughout, $x$ and $y$ denote complete
assignments, $x_i \in D_i$ is the value $x$ assigns to $X_i$, and
$\phi_i : D_i \to \mathbb{R}$ denotes the small unary tie-breaking
preferences used in the experiments (the dynamic unary constraint
$\phi^t$ of Section~\ref{sec:dynamic} is the time-varying case).  Scalar
statements (regimes, transition intervals) are for binary domains; claims
about larger domains are labeled empirical.

\begin{remark}[Schedule and orientation]\label{rem:conv}
We index messages by the two-phase schedule of
Section~\ref{secmaxsum}: in iteration $t$ all variable-nodes send
$Q^{t}$ (computed from $R^{t-1}$), then all function-nodes send $R^{t}$
(computed from $Q^{t}$); value selection at iteration $t$ uses $R^{t}$.
For cost terms, $C_{ij}$ denotes the \emph{shared} cost term of the
constraint between $X_i$ and $X_j$, read with $X_i$'s value first:
$C_{ij}(v, w) = C_{ji}(w, v)$ by definition.  No numerical symmetry of
the table is assumed.
\end{remark}

\subsection{From Regimes to Synchronous Local Selection}\label{sec:rule}

By Theorem~\ref{thm:stability}, outside the transition interval the
minimum in the $R$-update is attained at the same value of the sending
variable for both values of the receiving variable, so the $R$-message
equals one row of the cost table plus a constant.  For general domains we
say $Q^{t}_{X_i \to F}$ has a \emph{stable active minimizer} $u^{*}$ if
$u^{*} = \operatorname{argmin}_{u}[\,C'(u,v) + (Q^{t}_{X_i\to F})[u]\,]$
is the same for every $v$; then $R^{t}_{F\to X_j}(\cdot) =
C'(u^{*},\cdot) + (Q^{t})[u^{*}]$: \emph{the function-node forwards row
$u^{*}$}, up to an additive constant that shifts all entries equally
(this constant keeps propagating, which matters for damping).

\begin{lemma}[Clone synchronization]\label{lem:clones}
Under the symmetric split with zero-initialized messages, the two copies
$F'_{ij}, F''_{ij}$ send identical messages at every iteration, and so
does $X_i$ toward the two copies.
\end{lemma}

Lemma~\ref{lem:clones} exposes what splitting changes: the message $X_i$
sends toward $F'_{ij}$ \emph{contains the sibling copy's previous
$R$-message}.  Ordinary Min-sum is non-backtracking, whereas on the SCFG
every edge carries a length-two feedback cycle through the sibling copy.
A round trip through this cycle applies the same cost term twice, so it
reinforces its own direction whether the constraint prefers agreement or
disagreement, and its drift per round trip is the constant $d$ of
Section~\ref{sec:convergence}; Lemma~\ref{lem:drift} and
Theorem~\ref{thm:dyncvg} state that the cycle expels $\Delta Q$ from the
transition interval and holds it outside when the external swing obeys
the stated bound.  On benchmark graphs commitment is partial and
family-dependent: measured mean fractions of messages with a stable
active minimizer at iteration $9$ / in the tail are $0.73$/$0.82$
(random sparse, $|D_i|{=}10$, split) versus $0.12$/$0.10$ (unsplit),
$0.74$/$0.92$ versus $0.36$/$0.35$ (random dense), $0.97$/$0.97$ versus
$0.87$/$0.85$ (random binary), and $0.03$/$0.36$ versus $0.00$/$0.00$
(graph coloring, which rarely commits in this strict sense although its
messages still carry the sender's selection; Appendix~\ref{app:tworoutes}).
Commitment can also hold on one parity of an alternation and fail on the
other, and for $|D_i| > 2$ the commitment region is not an interval.

\begin{lemma}[Decoded assignment under full commitment]\label{lem:decode}
If at iteration $t$, for every neighbor $X_j$ of $X_i$, the message
$Q^{t}_{X_j \to F'_{ij}}$ has a stable active minimizer $u^{*t}_j$, then
\[
  \hat{x}^{t}_i \;=\; \operatorname{argmin}_{v \in D_i}
  \Big[\phi_i(v) + \sum_{j : C_{ij}\in\mathcal{C}}
  C_{ij}\big(v,\, u^{*t}_j\big)\Big].
\]
\end{lemma}

\begin{lemma}[Sibling evaluation]\label{lem:sibling}
Assume the hypothesis of Lemma~\ref{lem:decode} at iteration $t$, and
additionally that $Q^{t+1}_{X_i \to F'_{ij}}$ has a stable active
minimizer $u^{*,t+1}_i$, with unique argmins.  Then
$u^{*,t+1}_i = \hat{x}^{t}_i$.
\end{lemma}
\begin{proof}
By Lemma~\ref{lem:clones} and the hypothesis at $t$,
$(Q^{t+1}_{X_i \to F'_{ij}})[u] = \phi_i(u) + \sum_{k \neq j}
C_{ik}(u, u^{*t}_k) + \tfrac12 C_{ij}(u, u^{*t}_j) + \text{const}$, the
last term being the sibling's forwarded half-row.  A stable active
minimizer minimizes $\tfrac12 C_{ij}(u,v) + (Q^{t+1})[u]$ for
\emph{every} $v$; evaluating at $v = u^{*t}_j$ makes the two half-costs
of $C_{ij}$ recombine, so the minimized expression becomes the full field
$\phi_i(u) + \sum_{k} C_{ik}(u, u^{*t}_k) + \text{const}$, whose unique
argmin is $\hat{x}^{t}_i$ by Lemma~\ref{lem:decode}. \qedhere
\end{proof}

\begin{corollary}[The selection rule]\label{cor:rule}
If the hypotheses of Lemmas~\ref{lem:decode}--\ref{lem:sibling} hold for
every variable and constraint at every $t \ge t_0$ (\emph{full
commitment}), then for all $t \ge t_0$, simultaneously for all $i$,
\[
  \hat{x}^{t+1}_i \;=\; \operatorname{argmin}_{v\in D_i}
  \Big[\phi_i(v) + \sum_{j : C_{ij}\in\mathcal{C}}
  C_{ij}\big(v, \hat{x}^{t}_j\big)\Big]
\]
--- \emph{synchronous local selection}.
\end{corollary}

Because full commitment is partial in practice, we state the attractor
theorem under the weaker hypothesis that the selected assignments satisfy
this recursion --- implied by full commitment, and checkable directly on
a run.  On an audited subset of $32$ period-2 runs across four benchmark
families, the recursion held on every check performed: the two
alternating solutions mapped to one another under it exactly (both
directions, on the limit pair of every instance), and it predicted
$\hat{x}^{t+1}$ from $\hat{x}^{t}$ correctly on all sampled transitions
(every seventh transition of the recorded tails; protocol in
Appendix~\ref{app:tworoutes}).

\subsection{The Two-Solution Theorem}\label{sec:twosol}

\begin{definition}[Alternation cost]\label{def:altcost}
For complete assignments $(x, y)$ let
\[
  \mathrm{cost}_2(x,y) = \!\!\sum_{C_{ij}\in\mathcal{C}}\!\!
  \big[C_{ij}(x_i, y_j) + C_{ij}(y_i, x_j)\big]
  + \sum_i \big[\phi_i(x_i) + \phi_i(y_i)\big].
\]
Then $\mathrm{cost}_2(x,y) = \mathrm{cost}_2(y,x)$ --- using only that
the two endpoints share one cost term --- and $\mathrm{cost}_2(x,x) =
2\,\mathrm{cost}(x)$.
\end{definition}

Under the synchronous schedule a constraint is never evaluated within one
iteration: each variable reacts to its neighbors' \emph{previous} values,
so every charge crosses consecutive iterations.  $\mathrm{cost}_2(x,y)$ is
exactly the resulting bill for alternating between $x$ and $y$ --- and
since $\mathrm{cost}_2(x,x) = 2\,\mathrm{cost}(x)$, alternating pairs and
ordinary solutions are priced in the same currency, with diagonal pairs
being ordinary solutions counted twice.

\begin{theorem}[Two solutions]\label{thm:tworoutes}
If for all $t \ge t_0$ the selected assignments satisfy the synchronous
local selection recursion with unique argmins, then
$\mathrm{cost}_2(\hat{x}^{t}, \hat{x}^{t+1})$ is non-increasing and
strictly decreasing unless $\hat{x}^{t+2} = \hat{x}^{t}$; hence the
selected-assignment sequence is eventually periodic with period $1$ or
$2$.
\end{theorem}
\begin{proof}
Grouping $\mathrm{cost}_2(\hat{x}^{t+1}, z)$ around $z$, each $X_i$
contributes $\phi_i(z_i) + \sum_{j} C_{ij}(z_i, \hat{x}^{t+1}_j)$ plus
terms independent of $z$; this is separable and coordinatewise minimized
by $z = \hat{x}^{t+2}$.  Hence
$\mathrm{cost}_2(\hat{x}^{t+1}, \hat{x}^{t+2}) \le
\mathrm{cost}_2(\hat{x}^{t+1}, \hat{x}^{t}) =
\mathrm{cost}_2(\hat{x}^{t}, \hat{x}^{t+1})$ by swap symmetry.  With
unique argmins, equality forces $\hat{x}^{t+2} = \hat{x}^{t}$;
$\mathrm{cost}_2$ takes finitely many values, so strict decreases
terminate. \qedhere
\end{proof}

\begin{corollary}\label{cor:localmin}
The limit pair $(x, y)$ is a local minimum of $\mathrm{cost}_2$ under
single-variable changes in either coordinate.  Period $1$ occurs exactly
when the pair is diagonal, and a diagonal limit is a local minimum of the
original problem.
\end{corollary}

Parallel-update period-at-most-two results go back to
\citet{GolesOlivos80} (threshold networks, symmetric weights) and
\citet{PoljakSura83} (weighted-plurality opinion dynamics --- pairwise
terms proportional to equality indicators); a two-cycle result for
simultaneous best response in two-player random potential games appears
in \citet{AshkenaziGolan25}.  Theorem~\ref{thm:tworoutes} covers
arbitrary \emph{shared} pairwise cost terms with unaries, hence the short
proof.  Within its hypothesis it explains why the observed period is two
rather than three, five, or aperiodic.  Empirically (census in
Appendix~\ref{app:tworoutes}; periods detected on the final $150$ of
$700$ iterations, up to $64$): of $276$ split runs, $274$ had a detected
periodic tail of period $1$ or $2$ ($191$ period-2, $83$ period-1) and
two had longer detected periods ($16$, $42$); restricted to the
$n \ge 40$ cells of the domain-10 and coloring families the period-2
count is $92$ of $92$; unsplit runs on the random domain-10 families had
\emph{no} detected period at all in $85$--$100\%$ of instances from
$n{=}20$ up.

Off-diagonal minima of $\mathrm{cost}_2$ are where alternation pays: for
three mutually ``not-equal'' constrained variables, every single
assignment must violate one constraint, but the pair
$x = (\alpha,\alpha,\alpha)$, $y = (\beta,\beta,\beta)$ satisfies all six
cross terms.  On the audited subset, the limit pair was a local minimum
of $\mathrm{cost}_2$ with zero violating moves in $32/32$ runs.

\begin{corollary}[Bipartite rephasing]\label{cor:bipartite}
Let the constraint graph be connected and bipartite with sides $A, B$,
and $(x,y)$ any pair.  Define $x'$ ($=x$ on $A$, $=y$ on $B$) and $y'$
($=y$ on $A$, $=x$ on $B$).  Then
$\mathrm{cost}_2(x, y) = \mathrm{cost}(x') + \mathrm{cost}(y')$, and if
$(x,y)$ is a layer-wise local minimum of $\mathrm{cost}_2$ then $x'$ and
$y'$ are each coordinatewise local minima of the \emph{original} problem.
\end{corollary}

On a bipartite graph, then, the oscillating algorithm is not wandering
between two poor assignments: it holds \emph{two locally optimal
solutions in interleaved form}, and the raw snapshots read them out in
mixed phase.  This is exactly Observation~\ref{obs:lemniscate}
($\{X_1, X_3\}$ and $\{X_2\}$ are the two sides; snapshots cost $300$ and
$1200$ while their rephasings are the local optima of costs $130$ and
$132$), and it was verified on $8/8$ oscillating complete-bipartite
random instances (Appendix~\ref{app:tworoutes}).  On non-bipartite graphs
the only exact symmetry is the global swap $(x,y)\mapsto(y,x)$;
per-variable rephasing changes $\mathrm{cost}_2$ and is heuristic
post-processing (measured gains in Appendix~\ref{app:tworoutes}).  The
two-value menus of the merge heuristics of Section~\ref{Sec:MS-split} are
exactly the coordinates of the limit pair: on bipartite graphs
Corollary~\ref{cor:bipartite} guarantees they contain two coordinatewise
local optima; on general graphs the support is empirical
(Table~\ref{tab:merge}).

\section{Damping and the Two-Solution Oscillation}\label{sec:whydamping}

Damping is applied to the $Q$-messages
(Section~\ref{secmaxsum}).  Because a committed function-node forwards
its row plus a propagating constant, raw messages need not settle even
when all decisions are locked; the objects that settle are the message
\emph{differences}.  All statements below are about differences.

\begin{lemma}[Averaging an alternating computation]\label{lem:ema}
If a damped scalar sequence has computed values alternating between
constants $u$ and $v$, the sent values converge geometrically (rate
$\lambda$) to the alternating pair $\frac{u+\lambda v}{1+\lambda}$,
$\frac{v+\lambda u}{1+\lambda}$: the mean is preserved and the swing is
multiplied by $(1-\lambda)/(1+\lambda)$ --- $1/19$ at $\lambda = 0.9$.
\end{lemma}

\begin{theorem}[Persistence threshold of a committed alternation]\label{thm:lambdastar}
Binary domains.  Consider a trajectory of $Q$-damped Min-sum on the SCFG
that, from some iteration on, reproduces a fixed period-2 pattern of
regimes in which every message difference is in a regime at every step.
For a message whose regime side flips each iteration, the computed
differences alternate between constants $\Delta_1 > \tau_U$ and
$\Delta_2 < \tau_L$ determined by the pattern, and the pattern can
persist only if
\[
  \lambda \;\le\;
  \min\!\left(
    \frac{\Delta_1 - \tau_U}{\tau_U - \Delta_2},\;
    \frac{\tau_L - \Delta_2}{\Delta_1 - \tau_L}
  \right)
\]
for every flipping message; conversely, under strict inequalities the
corresponding period-2 orbit of damped message differences exists and is
locally attracting (contraction $\lambda$ per step).  The bound is
strictly below $1$ whenever $\tau_U > \tau_L$ or the two margins differ.
Damped and undamped updates have identical fixed points.
\end{theorem}

What Theorem~\ref{thm:lambdastar} does \emph{not} say: above the
threshold of a given pattern the trajectory need not converge ---
message differences may enter transition intervals, settle into a
different pattern, or sustain non-committed oscillations.  Indeed, of
$17$ initially oscillating instances in our damping scans, $2$ (both
random dense) were non-monotone in $\lambda$.  We therefore report per
instance the smallest grid $\lambda$ from which all larger tested values
produced a constant \emph{assignment} tail (final $300$ of $2000$
iterations; assignment constancy, cost constancy, and convergence of
message differences are increasingly strong statements and only the
first was tested): at most $0.2$ on random sparse and binary, with the
largest values, $0.4$, on three coloring instances and one random dense
instance (Appendix~\ref{app:tworoutes}).  For $|D_i| > 2$ the theorem
does not apply and the domain-10 results are empirical.

\begin{proposition}[The threshold approaches $1$ with degree]\label{prop:kn}
Let the constraint graph be complete on $n$ agents with binary domains,
every constraint costing $10$ when the two agents take equal values and
$0$ otherwise, symmetric split, and a uniform unary preference gap
$0.01$.  By symmetry all message differences share one value $q_t$,
whose damped dynamics is the scalar recurrence
\[
  q_{t+1} \;=\; \lambda q_t + (1-\lambda)\big[\,0.01 -
  (2n-3)\operatorname{clip}(q_t, -5, 5)\,\big],
\]
and the alternation between the two uniform solutions is a fully
committed period-2 pattern with persistence threshold
\[
  \frac{10(n-2) - 0.01}{10(n-1) + 0.01} \;\xrightarrow[n\to\infty]{}\; 1 .
\]
Hence for every fixed $\lambda < 1$ there is a degree at which the
oscillation persists: no problem-independent damping value exists.  For
$n = 12$ the threshold is $0.9089$; the customary $\lambda = 0.9$ does
not eliminate this oscillation (verified in the implementation from zero
initialization), while at $\lambda = 0.95$ the message differences
converge.
\end{proposition}

\paragraph{Summary, with scope.}
Under full commitment and clone synchronization, split Min-sum's
selections provably follow synchronous local selection
(Corollary~\ref{cor:rule}); any selection-following tail is eventually of
period $1$ or $2$ and parks at a local minimum of $\mathrm{cost}_2$
(Theorem~\ref{thm:tworoutes}), whose bipartite rephasings are ordinary
local optima (Corollary~\ref{cor:bipartite}).  On the benchmarks,
commitment is partial and family-dependent, yet the selection recursion
held on every check of the audited runs, and detected split tails were
overwhelmingly of period $1$ or $2$ --- period $2$ almost exclusively in
the large domain-10 and coloring cells --- while unsplit runs mostly
showed no detectable period.  For damping: alternating computed
differences are attenuated by $(1-\lambda)/(1+\lambda)$
(Lemma~\ref{lem:ema}); each fully committed binary period-2 pattern
persists exactly up to an explicit threshold
(Theorem~\ref{thm:lambdastar}); and the threshold approaches $1$ with
degree (Proposition~\ref{prop:kn}), so the reliability of
$\lambda = 0.9$ on the benchmark families --- measured stable-damping
values at most $0.4$ --- is a property of those families, not a universal
guarantee.  Damping removes the parallel-update alternation that
committed dynamics admits as its principal failure mode; it cannot
manufacture commitment, which is splitting's contribution --- partial,
but decisive on these benchmarks.
